\documentclass[11pt,a4paper]{article}
\usepackage[T2A]{fontenc}
\usepackage[utf8]{inputenc}
\usepackage[english,russian]{babel}
\sloppy
\setlength{\emergencystretch}{2em}
\usepackage[margin=2.1cm]{geometry}
\usepackage{xcolor}
\usepackage{graphicx}
\usepackage{listings}
\usepackage{enumitem}
\usepackage{booktabs}
\usepackage{array}
\newcolumntype{L}[1]{>{\raggedright\arraybackslash}p{#1}}
\usepackage{amsmath}
\usepackage{tikz}
\usetikzlibrary{arrows.meta,positioning,calc,fit,shapes.geometric}
\usepackage{titlesec}
\usepackage{fancyhdr}
\usepackage[hidelinks]{hyperref}
\usepackage{caption}
\captionsetup{font=small,labelfont=bf}
% babel нь \begin{document}-д эдгээрийг дардаг тул AtBeginDocument ашиглана
\AtBeginDocument{%
  \renewcommand{\contentsname}{Гарчиг}%
  \renewcommand{\tablename}{Хүснэгт}%
  \renewcommand{\figurename}{Зураг}%
  \renewcommand{\refname}{Ашигласан материал}%
}

\definecolor{codebg}{RGB}{246,247,249}
\definecolor{accent}{RGB}{31,78,121}

\lstset{
  basicstyle=\ttfamily\footnotesize,
  backgroundcolor=\color{codebg},
  frame=single, rulecolor=\color{gray!40},
  keywordstyle=\color{accent}\bfseries,
  commentstyle=\color{gray}\itshape,
  breaklines=true, columns=fullflexible, keepspaces=true,
  extendedchars=true, inputencoding=utf8,
  literate=
    {А}{{\CYRA}}1
    {Б}{{\CYRB}}1
    {В}{{\CYRV}}1
    {Г}{{\CYRG}}1
    {Д}{{\CYRD}}1
    {Е}{{\CYRE}}1
    {Ё}{{\CYRYO}}1
    {Ж}{{\CYRZH}}1
    {З}{{\CYRZ}}1
    {И}{{\CYRI}}1
    {Й}{{\CYRISHRT}}1
    {К}{{\CYRK}}1
    {Л}{{\CYRL}}1
    {М}{{\CYRM}}1
    {Н}{{\CYRN}}1
    {О}{{\CYRO}}1
    {Ө}{{\CYROTLD}}1
    {П}{{\CYRP}}1
    {Р}{{\CYRR}}1
    {С}{{\CYRS}}1
    {Т}{{\CYRT}}1
    {У}{{\CYRU}}1
    {Ү}{{\CYRY}}1
    {Ф}{{\CYRF}}1
    {Х}{{\CYRH}}1
    {Ц}{{\CYRC}}1
    {Ч}{{\CYRCH}}1
    {Ш}{{\CYRSH}}1
    {Щ}{{\CYRSHCH}}1
    {Ъ}{{\CYRHRDSN}}1
    {Ы}{{\CYRERY}}1
    {Ь}{{\CYRSFTSN}}1
    {Э}{{\CYREREV}}1
    {Ю}{{\CYRYU}}1
    {Я}{{\CYRYA}}1
    {а}{{\cyra}}1
    {б}{{\cyrb}}1
    {в}{{\cyrv}}1
    {г}{{\cyrg}}1
    {д}{{\cyrd}}1
    {е}{{\cyre}}1
    {ё}{{\cyryo}}1
    {ж}{{\cyrzh}}1
    {з}{{\cyrz}}1
    {и}{{\cyri}}1
    {й}{{\cyrishrt}}1
    {к}{{\cyrk}}1
    {л}{{\cyrl}}1
    {м}{{\cyrm}}1
    {н}{{\cyrn}}1
    {о}{{\cyro}}1
    {ө}{{\cyrotld}}1
    {п}{{\cyrp}}1
    {р}{{\cyrr}}1
    {с}{{\cyrs}}1
    {т}{{\cyrt}}1
    {у}{{\cyru}}1
    {ү}{{\cyry}}1
    {ф}{{\cyrf}}1
    {х}{{\cyrh}}1
    {ц}{{\cyrc}}1
    {ч}{{\cyrch}}1
    {ш}{{\cyrsh}}1
    {щ}{{\cyrshch}}1
    {ъ}{{\cyrhrdsn}}1
    {ы}{{\cyrery}}1
    {ь}{{\cyrsftsn}}1
    {э}{{\cyrerev}}1
    {ю}{{\cyryu}}1
    {я}{{\cyrya}}1,
  xleftmargin=0.8em, framexleftmargin=0.8em,
  language=Java, showstringspaces=false
}
\titleformat{\section}{\Large\bfseries\color{accent}}{\thesection}{0.6em}{}
\titleformat{\subsection}{\large\bfseries}{\thesubsection}{0.6em}{}
\pagestyle{fancy}
\fancyhf{}
\fancyfoot[C]{\thepage}
\renewcommand{\headrulewidth}{0.4pt}
\setlength{\headheight}{14pt}
\setlength{\parskip}{3pt}
\fancyhead[L]{\small IV-2-7 $\to$ IV-2-10}
\fancyhead[R]{\small Review}

\begin{document}
\begin{titlepage}
\centering
\vspace*{2.5cm}
{\large F.CSM302 --- Үйлдлийн систем · Бие даалт 1\par}
\vspace{1.4cm}
{\Huge\bfseries\color{accent} Харилцан шүүмж\par}
\vspace{1.0cm}
{\large\begin{tabular}{rl}
Шүүмжлэгч код: & \textbf{IV-2-7} \\[4pt]
Шүүмжилсэн код: & \textbf{IV-2-10} \\[4pt]
Төсөл: & The Sleeping Teaching Assistant \\[4pt]
 & (Бүлэг 7, UI дүрслэл: Queue animation) \\[4pt]
Огноо: & \rule{4cm}{0.4pt} \\
\end{tabular}}
\vfill
{\large Шүүмжлэгч багийн гишүүд\par}
\vspace{0.3cm}
\begin{tabular}{ll}
\toprule
Оюутны код & Нэр \\
\midrule
B200910032 & Д.Дагиймаа \\
B221940034 & Д.Оюунжаргал \\
B231930001 & \\
B241940014 & Г.Мягмарсүрэн \\
B241940127 & О.Энх \\
B242270084 & Э.Бүрэнбаяр \\
\bottomrule
\end{tabular}
\vspace{2cm}
\end{titlepage}

\tableofcontents
\newpage

%============================================================
\section{Ажлын ойлголт}
\emph{Зорилго ба хэрэгжилтийг ӨӨРИЙН үгээр 5--8 өгүүлбэрээр бич.
Тайланг нь хуулж болохгүй --- ойлгосноо харуул.}

\textbf{Бодлогын мөн чанар.} Sleeping Teaching Assistant нь сонгодог
Sleeping Barber бодлогын хувилбар: нэг туслах багш оффисдоо унтаж байдаг;
оюутан ирэхэд сэрж тусалдаг; завгүй бол оюутан коридорын гурван сандлын
нэгэнд суудаг; сандал дүүрэн бол оюутан явж дараа эргэж ирдэг.

\textbf{Синхрончлолын асуудал:}\\[3mm]
\noindent\rule{\textwidth}{0.4pt}

\vspace{2mm}\noindent\rule{\textwidth}{0.4pt}

\textbf{Багийн шийдэл:}\\[3mm]
\noindent\rule{\textwidth}{0.4pt}

\vspace{2mm}\noindent\rule{\textwidth}{0.4pt}

\vspace{2mm}\noindent\rule{\textwidth}{0.4pt}

%============================================================
\section{Туршилтын орчин}
\begin{center}
\small
\begin{tabular}{lL{9.5cm}}
\toprule
Үзүүлэлт & Утга \\
\midrule
OS & \\
Architecture & \\
Compiler / runtime & \\
Хувилбар & \\
Dependency & \\
Цөмийн тоо & \\
\bottomrule
\end{tabular}
\end{center}

%============================================================
\section{Build ба Run}
\emph{Manual-ийн алхмыг ӨӨРЧЛӨХГҮЙГЭЭР эхний оролдлогыг хий. Ажиллахгүй бол
тэр нь өөрөө олдвор (finding).}

\begin{center}
\small
\begin{tabular}{L{2.0cm}L{5.6cm}L{3.4cm}L{3.4cm}}
\toprule
Алхам & Команд & Хүлээгдсэн үр дүн & Бодит үр дүн \\
\midrule
Clone & & & \\[6pt]
Build & & & \\[6pt]
Run & & & \\[6pt]
\bottomrule
\end{tabular}
\end{center}

\textbf{Manual-аас гадуур хийх шаардлагатай болсон өөрчлөлт:}\\[3mm]
\noindent\rule{\textwidth}{0.4pt}

%============================================================
\section{Туршилт}

\subsection{Үндсэн тест (тайланд бичсэнийг давтах)}
\begin{center}
\small
\begin{tabular}{L{0.6cm}L{4.4cm}L{4.4cm}L{4.4cm}}
\toprule
\# & Оролт / тохиргоо & Хүлээгдсэн үр дүн & Бодит үр дүн \\
\midrule
1 & & & \\[6pt]
2 & & & \\[6pt]
3 & & & \\[6pt]
\bottomrule
\end{tabular}
\end{center}

\subsection{Нэмэлт тест (дор хаяж НЭГ, өөрсдөө сонгосон)}
\begin{center}
\small
\begin{tabular}{L{0.6cm}L{3.6cm}L{4.4cm}L{2.6cm}L{2.6cm}}
\toprule
\# & Оролт / тохиргоо & Яагаад сонгосон & Хүлээгдсэн & Бодит \\
\midrule
1 & сандал $= 0$ & Захын утга --- оюутан хэзээ ч суухгүй, шууд явах ёстой & & \\[6pt]
2 & сандал $< 0$ & Хүчингүй оролт --- ойлгомжтой алдаа өгөх ёстой & & \\[6pt]
3 & 50 оюутан зэрэг & Ачаалал --- дараалал халих эсэх & & \\[6pt]
4 & үйлчилгээний хугацаа $= 0$ & Хил --- хуваарилалт гацах эсэх & & \\[6pt]
\bottomrule
\end{tabular}
\end{center}

%============================================================
\section{Дүрслэлийн үнэлгээ}
\emph{Багшийн оноосон дүрслэл: \textbf{Queue animation}.}

\begin{center}
\small
\begin{tabular}{L{9.4cm}L{2.0cm}L{3.0cm}}
\toprule
Шалгах асуулт & Тийм/Үгүй & Нотолгоо \\
\midrule
Оролт, тохиргоог өөрчлөхөд дүрслэл шинэчлэгдэж байна уу? & & \\[5pt]
Дараалал дахь оюутны тоо бодит төлөвтэй нийцэж байна уу? & & \\[5pt]
TA-ийн төлөв (унтаж / тусалж) зөв харагдаж байна уу? & & \\[5pt]
Тэмдэглэгээ, legend, нэгж ойлгомжтой юу? & & \\[5pt]
Алхамчилж харах эсвэл эргэн шалгах боломжтой юу? & & \\[5pt]
Дүрслэл нь алгоритмыг тайлбарлах \emph{нотолгоо} болж байна уу? & & \\
\bottomrule
\end{tabular}
\end{center}

\textbf{Дүгнэлт:}\\[3mm]
\noindent\rule{\textwidth}{0.4pt}

%============================================================
\section{Давуу тал}
\emph{Дор хаяж ХОЁР, тус бүр НОТОЛГООТОЙ. ``Сайхан болсон'' гэж бичихгүй ---
юу, хаана, яаж гэдгийг зааж бич.}

\textbf{1.}\\[3mm]
\noindent\rule{\textwidth}{0.4pt}

\vspace{1mm}\emph{Нотолгоо:}\\[3mm]
\noindent\rule{\textwidth}{0.4pt}

\vspace{3mm}
\textbf{2.}\\[3mm]
\noindent\rule{\textwidth}{0.4pt}

\vspace{1mm}\emph{Нотолгоо:}\\[3mm]
\noindent\rule{\textwidth}{0.4pt}

%============================================================
\newpage
\section{Асуудал}
\emph{Асуудал бүрийг дараах дарааллаар бич:}
\begin{center}
Нөхцөл $\to$ Давтах алхам $\to$ Хүлээгдсэн үр дүн $\to$ Бодит үр дүн
$\to$ Нөлөө $\to$ Сайжруулах санал
\end{center}

\begin{center}
\small
\begin{tabular}{L{1.8cm}L{5.4cm}L{3.6cm}L{3.6cm}}
\toprule
Priority & Асуудал (давтах алхамтай) & Нөлөө & Сайжруулах санал \\
\midrule
Critical & & & \\[14pt]
Major & & & \\[14pt]
Minor & & & \\[14pt]
Suggestion & & & \\[14pt]
\bottomrule
\end{tabular}
\end{center}

\begin{center}
\small
\begin{tabular}{lL{11.0cm}}
\toprule
Түвшин & Утга \\
\midrule
Critical & Үндсэн функц ажиллахгүй эсвэл буруу үр дүн \\
Major & Гол тест, дүрслэл, Manual ноцтой дутуу \\
Minor & Үндсэн ажиллагаанд саад багатай чанарын асуудал \\
Suggestion & Нэмэлт сайжруулалтын санал \\
\bottomrule
\end{tabular}
\end{center}

%============================================================
\section{Сайжруулах санал}
\emph{Асуудал бүрт НЭГ хэрэгжүүлэх боломжтой санал. Удирдамжийн жишээ:}

\begin{quote}
\small
\textbf{Сул:} ``Дүрслэлээ сайжруул.'' --- нотолгоо, нөлөө, хийх ажил тодорхойгүй.

\textbf{Сайн:} ``Quantum $=0$ үед Run дарсны дараа програм давталтад орсон.
Quantum $\le 0$ үед Run-ийг идэвхгүй болгож, хүчинтэй утгын тайлбар
харуулах саналтай.''
\end{quote}

\textbf{1.}\\[3mm]
\noindent\rule{\textwidth}{0.4pt}

\vspace{3mm}\textbf{2.}\\[3mm]
\noindent\rule{\textwidth}{0.4pt}

\vspace{3mm}\textbf{3.}\\[3mm]
\noindent\rule{\textwidth}{0.4pt}

%============================================================
\section{Дүгнэлт}
\emph{Сайжруулах ГУРВАН гол зүйлийг эрэмбэлж бич.}

\textbf{1.}\\[3mm]
\noindent\rule{\textwidth}{0.4pt}

\vspace{3mm}\textbf{2.}\\[3mm]
\noindent\rule{\textwidth}{0.4pt}

\vspace{3mm}\textbf{3.}\\[3mm]
\noindent\rule{\textwidth}{0.4pt}

\vspace{6mm}
\noindent\textbf{Ерөнхий үнэлгээ:}\\[3mm]
\noindent\rule{\textwidth}{0.4pt}

%============================================================
\vfill
\noindent\rule{\textwidth}{0.4pt}\\[2pt]
\small
\textbf{Рубрик (өөрийгөө шалгах).}
Дахин ажиллуулах чадвар 20 · Зөв ажиллагааны шинжилгээ 20 ·
Туршилт ба UI дүрслэл 20 · Тайлан, Manual, repository 15 ·
Техникийн шүүмж 20 · Reply ба PDF бүрдэл 5.

\end{document}
